\documentclass[12pt,a4paper]{article}
\usepackage[utf8]{inputenc}
\usepackage[portuguese]{babel}
\usepackage{geometry}
\usepackage{graphicx}
\usepackage{listings}
\usepackage{xcolor}
\usepackage{hyperref}
\usepackage{fancyhdr}
\usepackage{titlesec}
\usepackage{enumitem}
\usepackage{amsmath}
\usepackage{amssymb}
\usepackage{float}
\usepackage{booktabs}
\usepackage{array}
\usepackage{longtable}
\UseRawInputEncoding
\setlength{\headheight}{14.5pt}

% Configurações de geometria
\geometry{margin=2.5cm}

% Configurações de cores para código
\definecolor{codegreen}{rgb}{0,0.6,0}
\definecolor{codegray}{rgb}{0.5,0.5,0.5}
\definecolor{codepurple}{rgb}{0.58,0,0.82}
\definecolor{backcolour}{rgb}{0.95,0.95,0.92}

% Configuração do lstlisting
\lstdefinestyle{mystyle}{
    backgroundcolor=\color{backcolour},   
    commentstyle=\color{codegreen},
    keywordstyle=\color{magenta},
    numberstyle=\tiny\color{codegray},
    stringstyle=\color{codepurple},
    basicstyle=\ttfamily\footnotesize,
    breakatwhitespace=false,         
    breaklines=true,                 
    captionpos=b,                    
    keepspaces=true,                 
    numbers=left,                    
    numbersep=5pt,                  
    showspaces=false,                
    showstringspaces=false,
    showtabs=false,                  
    tabsize=2
}

\lstset{style=mystyle}

% Configurações de cabeçalho e rodapé
\pagestyle{fancy}
\fancyhf{}
\rhead{Cadeira de Rodas Inteligente - PUC}
\lhead{Tutorial de Instalação e Uso}
\rfoot{Página \thepage}

% Configurações de títulos
\titleformat{\section}{\Large\bfseries}{\thesection}{1em}{}
\titleformat{\subsection}{\large\bfseries}{\thesubsection}{1em}{}

% Configurações do hyperref
\hypersetup{
    colorlinks=true,
    linkcolor=blue,
    filecolor=magenta,      
    urlcolor=cyan,
    pdftitle={Tutorial Cadeira de Rodas Inteligente},
    pdfpagemode=FullScreen,
}

\begin{document}

\begin{titlepage}
    \centering
    \vspace*{2cm}
    
    {\Huge\bfseries Cadeira de Rodas Inteligente\par}
    \vspace{0.5cm}
    
    {\Large\bfseries Tutorial Completo de Instalação e Uso\par}
    \vspace{1cm}
    
    {\large Projeto de Embarcados - PUC\par}
    \vspace{0.5cm}
    
    {\large Controle por Movimentos da Cabeça\par}
    \vspace{1cm}
    
    {\large Integração de sistemas de controle e visão computacional\\para automatização de cadeiras de rodas com comandos faciais\par}
    \vspace{1cm}

    
    \vfill
    
    {\large \today\par}
    
    \vspace{1cm}
    
    {\large \textbf{Autores:} Robson Duarte Vicente e Vinicius Oliveira Ribas\par}
    \vspace{0.5cm}
    
    {\large \textbf{Professor:} Mario Guimaraes Buratto\par}
    
\end{titlepage}

\tableofcontents
\newpage

\section{Introdução}

Este documento apresenta um tutorial completo para instalação, configuração e uso do sistema de controle de cadeira de rodas inteligente desenvolvido para a disciplina de Embarcados na PUC. O sistema utiliza visão computacional para detectar movimentos da cabeça do usuário e traduzi-los em comandos para uma cadeira de rodas simulada.

\subsection{Visão Geral do Sistema}

O projeto consiste em um sistema de controle diferencial que permite:

\begin{itemize}
    \item Detecção de movimentos da cabeça em tempo real
    \item Controle de direção (esquerda/direita) através de rotação da cabeça
    \item Controle de velocidade (frente/trás) através de inclinação da cabeça
    \item Interface gráfica intuitiva com feedback visual
    \item Integração com simulador CoppeliaSim
    \item Múltiplos modos de controle (automático e manual)
\end{itemize}

\subsection{Tecnologias Utilizadas}

\begin{table}[H]
\centering
\begin{tabular}{|l|l|}
\hline
\textbf{Tecnologia} & \textbf{Propósito} \\
\hline
Python 3.8+ & Linguagem principal \\
OpenCV & Processamento de imagem \\
MediaPipe & Detecção facial e pose \\
dlib & Detecção de landmarks faciais \\
Tkinter & Interface gráfica \\
CoppeliaSim & Simulação da cadeira de rodas \\
NumPy & Cálculos matemáticos \\
\hline
\end{tabular}
\caption{Tecnologias principais do projeto}
\end{table}

\section{Requisitos do Sistema}

\subsection{Requisitos de Hardware}

\begin{itemize}
    \item \textbf{Processador}: Intel i5 ou AMD equivalente (ou superior)
    \item \textbf{RAM}: Mínimo 8GB, recomendado 16GB
    \item \textbf{Webcam}: Resolução mínima 640x480, recomendado 720p ou superior
    \item \textbf{Espaço em disco}: 2GB livres para instalação
    \item \textbf{Sistema operacional}: Windows 10/11, Linux ou macOS
\end{itemize}

\subsection{Requisitos de Software}

\begin{itemize}
    \item \textbf{Python}: Versão 3.8 ou superior
    \item \textbf{CoppeliaSim}: Versão 4.5 ou superior (para simulação)
    \item \textbf{Webcam funcional} com drivers atualizados
    \item \textbf{Ambiente de desenvolvimento}: IDE ou editor de código
\end{itemize}

\section{Instalação}

\subsection{1. Preparação do Ambiente}

\subsubsection{Instalação do Python}

\begin{enumerate}
    \item Acesse \url{https://www.python.org/downloads/}
    \item Baixe a versão mais recente do Python 3.8+
    \item Durante a instalação, marque a opção "Add Python to PATH"
    \item Verifique a instalação executando no terminal:
\end{enumerate}

\begin{lstlisting}[language=bash]
python --version
pip --version
\end{lstlisting}

\subsubsection{Clone do Repositório}

\begin{lstlisting}[language=bash]
git clone https://github.com/duduw1/puc.git
cd embarcados-puc-cadeira-de-rodas
\end{lstlisting}

\subsection{2. Instalação das Dependências}

\subsubsection{Criação do Ambiente Virtual (Recomendado)}

\begin{lstlisting}[language=bash]
# Criar ambiente virtual
python -m venv venv

# Ativar ambiente virtual
# Windows:
venv\Scripts\activate
# Linux/macOS:
source venv/bin/activate
\end{lstlisting}

\subsubsection{Instalação dos Pacotes}

Crie um arquivo \texttt{requirements.txt} com o seguinte conteúdo:

\begin{lstlisting}
opencv-python==4.8.1.78
mediapipe==0.10.7
dlib==19.24.2
numpy==1.24.3
Pillow==10.0.1
pynput==1.7.6
tkinter
\end{lstlisting}

Execute a instalação:

\begin{lstlisting}[language=bash]
pip install -r requirements.txt
\end{lstlisting}

\subsection{3. Configuração de Arquivos Especiais}

\subsubsection{Download do Modelo dlib}

O arquivo \texttt{shape\_predictor\_68\_face\_landmarks.dat} é necessário para o funcionamento do dlib:

\begin{enumerate}
    \item Acesse: \url{http://dlib.net/files/shape_predictor_68_face_landmarks.dat.bz2}
    \item Baixe o arquivo (95MB)
    \item Extraia o arquivo .dat
    \item Coloque na pasta raiz do projeto
\end{enumerate}

\subsubsection{Configuração do CoppeliaSim}

\begin{enumerate}
    \item Baixe o CoppeliaSim em: \url{https://www.coppeliarobotics.com/downloads}
    \item Instale o software
    \item Abra o arquivo \texttt{FUNCIONAL COM TESTE 15.ttt} no CoppeliaSim
    \item Configure a porta 19997 para comunicação
\end{enumerate}

\section{Configuração do Sistema}

\subsection{1. Configuração da Webcam}

\subsubsection{Verificação da Webcam}

Execute o seguinte código para testar a webcam:

\begin{lstlisting}[language=python]
import cv2

# Testar webcam
cap = cv2.VideoCapture(0)
if cap.isOpened():
    print("Webcam funcionando corretamente")
    ret, frame = cap.read()
    if ret:
        print(f"Resolução: {frame.shape[1]}x{frame.shape[0]}")
    cap.release()
else:
    print("Erro ao acessar webcam")
\end{lstlisting}

\subsubsection{Ajuste de Configurações}

Se necessário, ajuste as configurações da webcam:

\begin{lstlisting}[language=python]
cap.set(cv2.CAP_PROP_FRAME_WIDTH, 640)
cap.set(cv2.CAP_PROP_FRAME_HEIGHT, 480)
cap.set(cv2.CAP_PROP_FPS, 30)
\end{lstlisting}

\subsection{2. Configuração do CoppeliaSim}

\subsubsection{Inicialização da Simulação}

\begin{enumerate}
    \item Abra o CoppeliaSim
    \item Carregue o arquivo \texttt{FUNCIONAL COM TESTE 15.ttt}
    \item Inicie a simulação (botão play)
    \item Verifique se a cadeira de rodas está visível na cena
\end{enumerate}

\subsubsection{Configuração da API}

O arquivo \texttt{sim.py} contém a API para comunicação com o CoppeliaSim. Verifique se a conexão está funcionando:

\begin{lstlisting}[language=python]
import sim
import simConst

# Testar conexão
sim.simxFinish(-1)
clientID = sim.simxStart('127.0.0.1', 19997, True, True, 5000, 5)
if clientID != -1:
    print("Conectado ao CoppeliaSim")
else:
    print("Falha na conexão")
\end{lstlisting}

\section{Modos de Operação}

\subsection{1. Modo Automático (MediaPipe)}

O arquivo \texttt{6-mediapipe\_melhorar.py} oferece o controle mais avançado:

\begin{lstlisting}[language=bash]
python 6-mediapipe_melhorar.py
\end{lstlisting}

\subsubsection{Características do Modo Automático}

\begin{itemize}
    \item Detecção facial com MediaPipe
    \item Controle de sensibilidade ajustável
    \item Interface gráfica completa
    \item Detecção de piscadas para ativação
    \item Sistema de trava por inatividade
    \item Feedback visual em tempo real
\end{itemize}

\subsubsection{Controles do Modo Automático}

\begin{table}[H]
\centering
\begin{tabular}{|l|l|}
\hline
\textbf{Movimento} & \textbf{Ação} \\
\hline
Inclinar cabeça para frente & Avançar \\
Inclinar cabeça para trás & Recuar \\
Girar cabeça para esquerda & Virar à esquerda \\
Girar cabeça para direita & Virar à direita \\
Piscar duas vezes & Ativar/desativar sistema \\
\hline
\end{tabular}
\caption{Controles do modo automático}
\end{table}

\subsection{2. Modo Manual (Mouse)}

O arquivo \texttt{1-move\_automatic.py} oferece controle por mouse:

\begin{lstlisting}[language=bash]
python 1-move_automatic.py
\end{lstlisting}

\subsubsection{Características do Modo Manual}

\begin{itemize}
    \item Controle por movimento do mouse
    \item Interface simples
    \item Sensibilidade ajustável
    \item Ideal para testes e calibração
\end{itemize}

\subsection{3. Modo dlib (Alternativo)}

O arquivo \texttt{5-dlib\_melhorado.py} usa dlib para detecção:

\begin{lstlisting}[language=bash]
python 5-dlib_melhorado.py
\end{lstlisting}

\subsubsection{Características do Modo dlib}

\begin{itemize}
    \item Detecção facial com dlib
    \item 68 pontos de referência facial
    \item Maior precisão, menor velocidade
    \item Requer arquivo de modelo específico
\end{itemize}

\section{Interface do Usuário}

\subsection{Elementos da Interface}

\subsubsection{Janela Principal}

A interface principal contém:

\begin{itemize}
    \item \textbf{Visualização da câmera}: Mostra o feed da webcam com landmarks
    \item \textbf{Setas direcionais}: Indicam a direção do movimento
    \item \textbf{Sliders de sensibilidade}: Ajustam a resposta do sistema
    \item \textbf{Indicadores de status}: Mostram o estado do sistema
\end{itemize}

\subsubsection{Controles de Sensibilidade}

\begin{itemize}
    \item \textbf{Sensibilidade Direita/Esquerda}: Controla a resposta aos movimentos horizontais
    \item \textbf{Sensibilidade Cima/Baixo}: Controla a resposta aos movimentos verticais
    \item \textbf{Faixa de valores}: 0.5 a 3.0 (padrão: 1.0)
\end{itemize}

\subsection{Indicadores Visuais}

\subsubsection{Cores das Setas}

\begin{itemize}
    \item \textbf{Vermelho claro}: Movimento suave
    \item \textbf{Vermelho médio}: Movimento moderado
    \item \textbf{Vermelho escuro}: Movimento intenso
\end{itemize}

\subsubsection{Status do Sistema}

\begin{itemize}
    \item \textbf{Verde}: Sistema ativo e funcionando
    \item \textbf{Vermelho}: Sistema bloqueado ou com erro
    \item \textbf{Amarelo}: Sistema em calibração
\end{itemize}

\section{Calibração e Ajustes}

\subsection{1. Calibração Inicial}

\subsubsection{Posicionamento}

\begin{enumerate}
    \item Sente-se a uma distância de 50-80cm da webcam
    \item Mantenha o rosto centralizado na tela
    \item Evite iluminação muito forte ou muito fraca
    \item Remova óculos refletivos se possível
\end{enumerate}

\subsubsection{Calibração do Sistema}

\begin{enumerate}
    \item Execute o programa
    \item Mantenha a cabeça em posição neutra por 3 segundos
    \item O sistema detectará automaticamente a posição de referência
    \item Aguarde a mensagem "Calibração concluída"
\end{enumerate}

\subsection{2. Ajustes de Sensibilidade}

\subsubsection{Ajuste Fino}

\begin{enumerate}
    \item Comece com sensibilidade 1.0
    \item Teste movimentos pequenos da cabeça
    \item Aumente a sensibilidade se necessário
    \item Diminua se os movimentos forem muito sensíveis
\end{enumerate}

\subsubsection{Valores Recomendados}

\begin{table}[H]
\centering
\begin{tabular}{|l|l|}
\hline
\textbf{Tipo de Usuário} & \textbf{Sensibilidade Recomendada} \\
\hline
Iniciante & 0.8 - 1.2 \\
Intermediário & 1.0 - 1.5 \\
Avançado & 1.5 - 2.0 \\
\hline
\end{tabular}
\caption{Sensibilidade por nível de experiência}
\end{table}

\section{Solução de Problemas}

\subsection{1. Problemas Comuns}

\subsubsection{Webcam não detectada}

\begin{lstlisting}[language=python]
# Verificar dispositivos disponíveis
import cv2

for i in range(10):
    cap = cv2.VideoCapture(i)
    if cap.isOpened():
        print(f"Webcam encontrada no índice {i}")
        cap.release()
\end{lstlisting}

\subsubsection{Erro de conexão com CoppeliaSim}

\begin{enumerate}
    \item Verifique se o CoppeliaSim está rodando
    \item Confirme a porta 19997 está livre
    \item Reinicie o CoppeliaSim
    \item Verifique o firewall do sistema
\end{enumerate}

\subsubsection{Baixa performance}

\begin{itemize}
    \item Reduza a resolução da webcam
    \item Feche outros programas
    \item Use o modo dlib em vez do MediaPipe
    \item Ajuste os parâmetros de detecção
\end{itemize}

\subsection{2. Logs e Debug}

\subsubsection{Habilitar Logs Detalhados}

Adicione ao código:

\begin{lstlisting}[language=python]
import logging

logging.basicConfig(level=logging.DEBUG)
logger = logging.getLogger(__name__)
\end{lstlisting}

\subsubsection{Verificação de Performance}

\begin{lstlisting}[language=python]
import time

start_time = time.time()
# Código a ser medido
end_time = time.time()
print(f"Tempo de execução: {end_time - start_time:.3f} segundos")
\end{lstlisting}

\section{Desenvolvimento e Extensões}

\subsection{1. Estrutura do Código}

\subsubsection{Arquivos Principais}

\begin{itemize}
    \item \texttt{6-mediapipe\_melhorar.py}: Versão mais avançada
    \item \texttt{5-dlib\_melhorado.py}: Versão com dlib
    \item \texttt{1-move\_automatic.py}: Controle por mouse
    \item \texttt{sim.py}: API do CoppeliaSim
    \item \texttt{simConst.py}: Constantes da API
\end{itemize}

\subsubsection{Funções Principais}

\begin{lstlisting}[language=python]
# Funções de detecção
def calculate_pitch_using_nose(landmarks)
def calculate_yaw_using_eyes(landmarks)
def eye_aspect_ratio(landmarks, eye_indices)

# Funções de controle
def send_coppelia_command(left_wheel_speed, right_wheel_speed)
def update_arrows_based_on_control(yaw_control, pitch_control)

# Funções de interface
def create_arrow(canvas, direction, color, intensity)
def update_camera_frame(frame)
\end{lstlisting}

\subsection{2. Personalizações}

\subsubsection{Modificar Controles}

Para alterar os controles, edite as funções de mapeamento:

\begin{lstlisting}[language=python]
def get_direction_names(yaw_control, pitch_control):
    # Personalizar aqui os nomes das direções
    pass
\end{lstlisting}

\subsubsection{Adicionar Novos Modos}

Crie novos arquivos seguindo a estrutura:

\begin{lstlisting}[language=python]
# Importações necessárias
import cv2
import numpy as np
import tkinter as tk
# ... outras importações

# Função principal
def main():
    # Inicialização
    # Loop principal
    # Interface
    pass

if __name__ == "__main__":
    main()
\end{lstlisting}

\section{Considerações de Segurança}

\subsection{1. Segurança Física}

\begin{itemize}
    \item Sempre teste em ambiente simulado primeiro
    \item Mantenha distância segura de obstáculos
    \item Tenha um método de parada de emergência
    \item Monitore o sistema constantemente
\end{itemize}

\subsection{2. Segurança de Dados}

\begin{itemize}
    \item Não compartilhe imagens da webcam
    \item Use o sistema apenas em ambientes privados
    \item Desative a câmera quando não estiver em uso
    \item Mantenha o software atualizado
\end{itemize}

\section{Suporte}

Para dúvidas e suporte:

\begin{itemize}
    \item \textbf{Issues}: Use o sistema de issues do GitHub
    \item \textbf{Documentação}: Consulte este tutorial
    \item \textbf{Comunidade}: Participe dos fóruns de discussão
\end{itemize}

\section{Autores}

\begin{itemize}
    \item \textbf{Robson Duarte Vicente} - Desenvolvimento e implementação
    \item \textbf{Vinicius Oliveira Ribas} - Desenvolvimento e implementação
    \item \textbf{Mario Guimaraes Buratto} - Orientação e supervisão
\end{itemize}

\section{Referências e Links}

\subsection{Repositório do Projeto}

O código fonte completo deste projeto está disponível no GitHub:

\begin{itemize}
    \item \textbf{Repositório Principal}: \url{https://github.com/duduw1/PUC/tree/main/embarcados-puc-cadeira-de-rodas}
    \item \textbf{Issues e Discussões}: Use o sistema de issues do GitHub para reportar problemas ou sugerir melhorias
    \item \textbf{Contribuições}: Pull requests são bem-vindos para melhorar o projeto
\end{itemize}

\subsection{Recursos Adicionais}

\begin{itemize}
    \item \textbf{CoppeliaSim}: \url{https://www.coppeliarobotics.com/}
    \item \textbf{MediaPipe}: \url{https://mediapipe.dev/}
    \item \textbf{dlib}: \url{http://dlib.net/}
    \item \textbf{OpenCV}: \url{https://opencv.org/}
\end{itemize}

\section{Conclusão}

Este tutorial apresentou uma visão completa do sistema de controle de cadeira de rodas inteligente. O projeto demonstra a aplicação prática de visão computacional e sistemas embarcados para melhorar a qualidade de vida de pessoas com mobilidade reduzida.

\subsection{Próximos Passos}

\begin{itemize}
    \item Teste todos os modos de operação
    \item Ajuste as configurações conforme necessário
    \item Explore as possibilidades de personalização
    \item Considere contribuir para o desenvolvimento
\end{itemize}

\appendix

\section{Apêndice A: Comandos Úteis}

\subsection{Comandos do Terminal}

\begin{lstlisting}[language=bash]
# Verificar versão do Python
python --version

# Listar pacotes instalados
pip list

# Atualizar pip
python -m pip install --upgrade pip

# Instalar pacote específico
pip install nome_do_pacote

# Criar requirements.txt
pip freeze > requirements.txt
\end{lstlisting}

\subsection{Comandos do CoppeliaSim}

\begin{itemize}
    \item \textbf{Play}: Iniciar simulação
    \item \textbf{Pause}: Pausar simulação
    \item \textbf{Stop}: Parar simulação
    \item \textbf{Reset}: Reiniciar simulação
\end{itemize}

\section{Apêndice B: Configurações Avançadas}

\subsection{Parâmetros de Detecção}

\begin{lstlisting}[language=python]
# Parâmetros MediaPipe
mp_pose.Pose(min_detection_confidence=0.7)
mp_face_mesh.FaceMesh(refine_landmarks=True)

# Parâmetros dlib
EAR_THRESHOLD = 0.20
CONSEC_FRAMES = 1
\end{lstlisting}

\subsection{Configurações de Performance}

\begin{lstlisting}[language=python]
# Reduzir resolução para melhor performance
cap.set(cv2.CAP_PROP_FRAME_WIDTH, 320)
cap.set(cv2.CAP_PROP_FRAME_HEIGHT, 240)

# Ajustar FPS
cap.set(cv2.CAP_PROP_FPS, 15)
\end{lstlisting}

\end{document} 